\subsection*{3.3 Bounded-Degree Polynomials}
Robere and Zuiddam [12,Theorem5.13] showed that if $f$ is a polynomial of degree $d<n$, it is possible to improve on Potechin's theorem by decreasing $m=2^{2^n-1}$ down to $m=2^{\binom n{\leq d}-1}$. Here we show how to adapt Algorithm1 to get a similar result, and then at the end of the section we build a tradeoff algorithm to improve it.
\paragraph{Theorem13.} For any function $f:\{0,1\}^n\to\{0,1\}$ which is a degree-$d$ polynomial, there is an $m$-catalytic branching program with length $4n$ and width $2m$ that computes $f$, where $m\leq2^{n+\binom n{\leq d}}$.
Again, while this is slightly worse than Robere and Zuiddam's original result, we include it to show the flexibility of our approach and as a stepping stone to our tradeoff result.
\begin{proof}
The proof can be summarized as follows. We start with Algorithm1, and make the following change: for every $S'\in[n]$ such that $c_{S,S'}=0$ for all $S$, remove the register $R_{S'}$. We then argue that only $n+\binom n{\leq d}$ registers remain.
Let $p_f=f\in\mathbb F_2[x_1,\ldots,x_n]$ – since $f$ is already a polynomial, there's no need to define $p_f$ using interpolation this time. As before, let
\[q_f(\vec\tau^{\mathrm{in}},\vec y)=p_f(\vec y-\vec\tau^{\mathrm{in}})=\sum_{S,S'\subseteq[n]}c_{S,S'}\left(\prod_{i\in S}\tau_i^{\mathrm{in}}\right)\left(\prod_{i\in S'}y_i\right).\]
Let $M_f\subseteq2^{[n]}$ be the set of all $S'$ such that there exists an $S$ where $c_{S,S'}\neq0$. We define the following subroutine:
\paragraph{ToggleUsefulMonomials.}
\begin{enumerate}[label=\arabic*:]
\item \textbf{procedure} ToggleUsefulMonomials
\item ToggleInput
\item \textbf{for} $S\in M_f$ \textbf{do}
\item $R_S\leftarrow R_S+\prod_{i\in S}R_i^{\mathrm{in}}$
\item \textbf{end for}
\item ToggleInput
\item \textbf{end procedure}
\end{enumerate}
The only difference from ToggleMonomials is that we ignore subsets $S$ which are not in $M_f$ (not ``useful''). Our final algorithm is
\paragraph{Algorithm4: Transparently compute a degree-$d$ polynomial $f$ in time $4n$ with $n+\binom n{\leq d}$ registers.}
\begin{enumerate}[label=\arabic*:]
\item ToggleUsefulMonomials
\item $R_{\mathrm{out}}\leftarrow R_{\mathrm{out}}+\sum_{S\subseteq[n]}\sum_{S'\in M_f}c_{S,S'}(\prod_{i\in S}R_i^{\mathrm{in}})R_{S'}$
\item ToggleUsefulMonomials
\item $R_{\mathrm{out}}\leftarrow R_{\mathrm{out}}-\sum_{S\subseteq[n]}\sum_{S'\in M_f}c_{S,S'}(\prod_{i\in S}R_i^{\mathrm{in}})R_{S'}$
\end{enumerate}
To conclude the proof of Theorem13, we need to show $|M_f|\leq\binom n{\leq d}$. Indeed, since $p_f$ is a degree-$d$ polynomial, $q_f$ also has degree $d$, which means $c_{S,S'}=0$ whenever $|S|+|S'|>d$. So, $M_f$ only contains sets $S'$ with size at most $d$, of which there are $\binom n{\leq d}$.
\end{proof}
\paragraph{Note14.} The original algorithms of [11,12] rely on the symmetries of $f$ as an $\mathbb F_2$ polynomial, in essence having each start state represent a possible function $g$ which can be obtained from $f$ by negating input variables or taking $\oplus$ with $f$ itself. [11] takes this set of functions to be the space of all $n$-variable functions, while [12] analyzes these rules and obtains a more exact characterization. While this characterization is phrased in terms of orbit closures, it can also be described in terms of polynomials as the span of the set of all monomials appearing in $f$ as an $\mathbb F_2$ polynomial along with all submonomials of this set; this exactly coincides with our notion as $\prod_{i\in S}y_i$ generates all submonomials $\prod_{i\in S'\subseteq S}x_i$ for $y_i:=x_i+\tau_i$, which leads to the quantitative results being essentially the same despite taking two completely different approaches.
Now we state our tradeoff algorithm, which goes much in the same way as Theorem11 but without breaking the variables into groups.
\begin{proof}[Proof of Theorem15]
For any $\Delta\in\mathbb N$, let $M_f^{\Delta}\subseteq\binom n{\leq\Delta}$ be the set of all $S''$ of size at most $\Delta$ such that there exists an $S\subseteq[n]$ and $S'\supseteq S''$ where $c_{S,S'}\neq0$. We will have $k$ registers $R_{j,S''}$ for every $S''\in M_f^{\lceil d/k\rceil}$, as well as the usual registers $R_1^{\mathrm{in}}\ldots R_n^{\mathrm{in}},R_{\mathrm{out}}$. Note that this gives us our target space, as $|M_f^{\lceil d/k\rceil}|\leq\binom n{\leq\lceil d/k\rceil}$.
Following our proof of Theorem11, let $z_{j,S}=\tau_{j,S}+\prod_{i\in S}y_i$, where $\tau_{j,S}$ is the initial value of register $R_{j,S}$, and for every monomial in $q_f$ we split the term $\prod_{i\in S'}y_i$ in the monomial arbitrarily into $k$ different products $\prod_{i\in S'_j}y_i$ – each of which we can replace with $z_{j,S'_j}-\tau_{j,S'_j}$ – where $S'_j\in M_f^{\lceil d/k\rceil}$ and $\cup_{j\in[k]}S'_j=S'$. This is possible because each non-zero term in $q_f$ has degree at most $d$, meaning that $|S'|\leq d$ and furthermore every subset of $S'$ of size at most $\lceil d/k\rceil$ appears in $M_f^{\lceil d/k\rceil}$ by construction.\footnote{Our use of $j$ here is slightly different than in our previous proof; namely, $j$ is not linked to a specific block of variables, and rather we arbitrarily partitioned $S'$ into $k$ sets and assigned them each a distinct $j$. This will result in us having to spend time $n$ to load the monomials in, rather than time $\lceil n/k\rceil$ as in the previous proof, but this is necessary as we have no guarantee that there is a partition of the variables such that every monomial of degree at most $d$ is split into $k$ monomials of degree at most $\lceil d/k\rceil$. Note that this is where our algorithm performs worse than Algorithm4 when $d=\Omega(n)$.}
Fixing some particular partition $(S'_j)_{j\in[k]}$ for each $S'$, this gives us a new polynomial
\begin{multline*}
r_f(\vec\tau^{\mathrm{in}},\vec\tau,\vec z)=\sum_{S\subseteq[n]}\sum_{S'\subseteq[n]}\sum_{T\subseteq[k]}d_{S,S',T}\\
\cdot\left(\prod_{i\in S}\tau_i^{\mathrm{in}}\right)\left(\prod_{j\in T}z_{j,S'_j}\right)\left(\prod_{j\in[k]\setminus T}\tau_{j,S'_j}\right)
\end{multline*}
which is equivalent to $f(x_1\ldots f_n)$ as long as $z_{j,S'_j}=\tau_{j,S'_j}+\prod_{i\in S'_j}(x_i+\tau_i^{\mathrm{in}}+1)$. We define ToggleMonomialsForGroup as before, using ToggleInput instead of ToggleInputForGroup since the variables are no longer split into groups, and using a Gray code we get our final algorithm:
\paragraph{Algorithm5: Transparently compute $f$ in time $2^k\cdot2n$ with $n+k\cdot\binom n{\leq\lceil d/k\rceil}$ registers.}
\begin{enumerate}[label=\arabic*:]
\item \textbf{for} $\ell=0,\ldots,2^k-1$ \textbf{do}
\item \[\begin{aligned}R_{\mathrm{out}}\leftarrow R_{\mathrm{out}}+\sum_{S\subseteq[n]}\sum_{S'\subseteq[n]}&d_{S,S',T_{\ell}}\\
&\mathord{}(\prod_{i\in S}R_i^{\mathrm{in}})(\prod_{j=1}^kR_{j,S'_j})\end{aligned}\]
\item ToggleMonomialsForGroup($e_{\ell}$)
\item \textbf{end for}
\end{enumerate}
The analysis is identical to that of Algorithm4, except the runtime is $2^k\cdot2n$ rather than $2^k\cdot2\lceil n/k\rceil$ because we do not split the variables into groups.
\end{proof}
